%%%%%%%%%%%%%%%%%%%%%%%%%%%%%%%%%%%%%%%%%%%%%%%%%%%%%%%%%%%%%%%%%%%%%%%%%%%%%%%%
%   Copyright 2009, Oracle and/or its affiliates.
%   All rights reserved.
%
%
%   Use is subject to license terms.
%
%   This distribution may include materials developed by third parties.
%
%%%%%%%%%%%%%%%%%%%%%%%%%%%%%%%%%%%%%%%%%%%%%%%%%%%%%%%%%%%%%%%%%%%%%%%%%%%%%%%%

\newpage
\section{\overloadingcore}
\applabel{overloading-calculus}

In this section, we define a Fortress core calculus with overloading
for dotted methods and first-order functions.
We call this calculus \emph{\overloadingcore}.
\overloadingcore\ is an extension of \basiccore\ with overloading.
\revision{revival-calculi}{This calculus predates instantiation exclusion
(\secref{trait-types}, \secref{trait-decls}).
Its only condition on multiple inheritance is its rule
\ruleLabel{ValidMeth}, which constrains pairs of visible methods of one
name, so, as in \basiccore, an object that extends two different
instantiations of a parameterized trait declaring no method is well typed,
and its subtyping rule \sBothRule\ makes the object a subtype of both
instantiations; under the rule it is not a valid Fortress program.
The calculus, its rules and its soundness claim are as the team wrote and
proved them, and are not changed.}

\renewcommand{\tappone}{\ensuremath{N}}
\renewcommand{\tapptwo}{\ensuremath{M}}
\renewcommand{\tappthree}{\ensuremath{L}}
\renewcommand{\tappfour}{\ensuremath{K}}
\input{\macros/overloading-calculus-macros}
\input{\calculi/overloading/syntax}
\input{\calculi/overloading/dynamic}
\input{\calculi/overloading/static}
